\chapter*{LỜI MỞ ĐẦU}
\addcontentsline{toc}{chapter}{LỜI MỞ ĐẦU}

\section*{Lý do chọn đề tài}

Trong hoạt động tiếp thị số hiện đại, một chiến dịch marketing thường trải qua nhiều công đoạn phức tạp và có tính liên đới cao: xác định mục tiêu kinh doanh, phân bổ ngân sách, lựa chọn kênh truyền thông, sáng tạo nội dung đa định dạng, phê duyệt phân tầng, lập lịch phát hành và liên tục giám sát các chỉ số hiệu quả chuyển đổi (views, clicks, conversions, cost, revenue). Khi dữ liệu bị phân tán trên nhiều bảng tính hoặc các nền tảng rời rạc, đội ngũ marketing gặp nhiều trở ngại trong việc đối chiếu phiên bản nội dung, bảo đảm tính toàn vẹn của số liệu và tổng hợp báo cáo đa chiều. Sự ra đời của các mô hình ngôn ngữ lớn (LLMs) mở ra tiềm năng tự động hóa mạnh mẽ, nhưng việc ứng dụng AI thiếu kiểm soát tiềm ẩn rủi ro lớn về hiện tượng ảo giác (hallucination), vi phạm bảo mật dữ liệu và thiếu cơ chế truy cứu trách nhiệm.

Đề tài \textbf{“Hệ thống quản lý chiến dịch marketing có tích hợp trí tuệ nhân tạo (MarketFlow AI)”} được triển khai nhằm giải quyết triệt để bài toán này. Hệ thống kết hợp kiến trúc quản trị dữ liệu chặt chẽ theo chuẩn RESTful với quy trình Human-in-the-loop (HITL) nghiêm ngặt, tích hợp hệ sinh thái mô hình Google Gemini để hỗ trợ người dùng ở ba tác vụ then chốt: gợi ý ý tưởng chiến dịch, tạo bản nháp nội dung chuẩn hóa theo kênh, và tóm tắt phân tích hiệu quả dựa trên chỉ số thực tế.

\section*{Mục đích và phạm vi của Báo cáo V9}

Báo cáo V9 là tài liệu học thuật hoàn chỉnh tổng kết toàn diện quá trình nghiên cứu, thiết kế, hiện thực hóa mã nguồn và đánh giá thực nghiệm của dự án MarketFlow AI. Hệ thống không dừng lại ở mức mô hình phân tích lý thuyết mà đã được lập trình hoàn chỉnh toàn bộ mã nguồn (Full-stack), đóng gói chuẩn container Docker và trải qua các vòng kiểm thử đối kháng khắt khe.

Tài liệu được cấu trúc chặt chẽ gồm 8 chương chuẩn mực:
\begin{enumerate}
  \item \textbf{Chương 1: Giới thiệu tổng quan đề tài và bối cảnh} --- Phân tích bối cảnh tiếp thị số, các bên liên quan, mục tiêu đo lường được và ranh giới hệ thống.
  \item \textbf{Chương 2: Cơ sở kỹ thuật và công nghệ} --- Khảo sát nền tảng công nghệ cốt lõi: FastAPI, React 18, hệ sinh thái Google Gemini LLM, SQLite (WAL mode) và Docker containerization.
  \item \textbf{Chương 3: Phân tích yêu cầu hệ thống} --- Đặc tả chi tiết 14 yêu cầu chức năng (FR01--FR14), 6 yêu cầu chất lượng (NFR01--NFR06), ma trận phân quyền RBAC kết hợp kiểm soát mức bản ghi (Record-level NFR01) và máy trạng thái duyệt nội dung HITL.
  \item \textbf{Chương 4: Thiết kế hệ thống} --- Trình bày kiến trúc đa tầng (C4 Container), mô hình dữ liệu quan hệ ERD, luồng dữ liệu nghiệp vụ và thiết kế Prompt Engine kiểm soát ngữ cảnh.
  \item \textbf{Chương 5: Hiện thực hóa hệ thống} --- Chi tiết cài đặt các module backend FastAPI, giao diện web React 18 / Tailwind CSS và cơ chế AI Smart Fallback khử ảo giác (De-hallucination).
  \item \textbf{Chương 6: Kiểm thử thực nghiệm} --- Báo cáo kết quả bộ kiểm thử tự động (Unit, Integration, Security, Adversarial) với 205 ca kiểm thử, tỷ lệ đạt 100\%, 0 lỗi và 0 cảnh báo.
  \item \textbf{Chương 7: Đánh giá AI định lượng} --- Phân tích số liệu đo đạc thực nghiệm: Schema validation rate (100\%), Grounding score (>95\%), độ trễ P50/P95 và chi phí thực tế 0 VNĐ trên Google AI Studio Free Tier.
  \item \textbf{Chương 8: Kết luận và giới hạn đã biết} --- Đánh giá khách quan thành quả đạt được, phân tích trung thực các điểm giới hạn kỹ thuật và định hướng nâng cấp trong tương lai.
\end{enumerate}

\section*{Đóng góp kỹ thuật và phương pháp tiếp cận thực chứng}

Dự án MarketFlow AI tuân thủ phương pháp nghiên cứu khoa học thiết kế (Design Science Research Methodology -- DSRM) kết hợp chuẩn kỹ nghệ phần mềm doanh nghiệp, mang lại bốn đóng góp kỹ thuật cốt lõi:
\begin{itemize}[noitemsep,topsep=2pt]
  \item \textbf{Kiểm soát truy cập mức bản ghi (Record-level Access Control)}: Giải quyết triệt để nguy cơ lộ lọt dữ liệu chéo giữa các chiến dịch (IDOR), bảo đảm nguyên tắc đặc quyền tối thiểu cho từng nhân viên tiếp thị.
  \item \textbf{Máy trạng thái duyệt nội dung bất biến (HITL State Machine)}: Đóng kín mọi ngả đường bypass trạng thái xuất bản, bảo đảm 100\% nội dung trước khi ra công chúng đều được người có thẩm quyền phê chuẩn.
  \item \textbf{Prompt Engineering khử ảo giác (De-hallucination Engine)}: Ràng buộc cấu trúc đầu ra nghiêm ngặt qua Pydantic v2 và thuật toán nội suy chỉ số thuần túy số học, loại bỏ các kết luận tự suy đoán vô căn cứ.
  \item \textbf{Tối ưu hóa kinh tế vận hành}: Khai thác hiệu quả hạn ngạch Google AI Studio Free Tier với chi phí 0 VNĐ, xây dựng mô hình dự phóng tài chính minh bạch cho bài toán mở rộng quy mô.
\end{itemize}

Toàn bộ các nhận định, sơ đồ kiến trúc và bảng số liệu trong báo cáo đều được đối chiếu trực tiếp với mã nguồn thực tế và kết quả đo đạc từ môi trường kiểm thử của dự án.
